\section{Related Work}
It is not uncommon to have ``intentional unsoundness'' in industrial languages,
such as Java's covariant array subtyping, or TypeScript's covariant method subtyping.
However these are intentionally designed into the language to improve ease of use, and in most cases (like Java's array handling)
problematic uses are detected at runtime (e.g., Java checks upon storing into an array that the new entry value
is a subtype of the array's original content type).

However, unintentional unsoundness in industrial type systems is rarer.
The most notable example is Amin and Tate's discovery~\cite{amin2016java} that Java and Scala's handling of generic type constraints
was rendered unsound by the fact that \Q@null@ is treated as inhabiting all reference types, making all
type constraints (even nonsensical ones) satisfiable.
Servetto et al.~\cite{unsound_presentation} first pointed out the issue publicly, though in the 
Q\&A session after the talk, multiple Java designers were unsurprised --- as Gilad Bracha put it, ``Once I saw the title of your talk, I understood what the problem was.''
\todo[inline]{Colin says: As I listened to the Q\&A session again, James framed this as a question of whether 
two instances allocated in different contexts should be considered the same or not. I think that might be useful framing to bring up early here.
Gilad also pointed out that (paraphrasing) the BETA people knew an inner class was a property of the instance, not its own class.
}

\todo[inline]{discuss C\# casts with reified generics}

Flix~\tocite{Flix FAQ, for monomorphization and no subtyping} is a functional language for the JVM.
The Flix compiler monomorphizes polymorphic functions, as Rust and MLton do.
Flix has no subtyping.
The Flix construct closest to a local class is the \Q@new@ expression, which creates an object of an unnamed class that extends a Java class or implements a Java interface~\tocite{Flix anonymous classes}.
A Flix program can not name such a class, so a Flix program can not write a type test against it, as with the generated classes of C\#.
Thus the type test of Example~1 can not be written in Flix.

The key issues at work in the interaction of method-local inner classes and generics arise from the interaction of (1) nominal dynamic downcasts, (2) design choices about what runtime information about generic parameters is both persisted \emph{and checked} by casts, and (3) (not) tracking all dependencies on type parameters in the type system.
Essentially, Java and Scala both treat \Q@MyFoo@ as a sort of existential type --- that there exists some \Q@X@  such that
\Q@MyFoo@ extends \Q@Foo<X>@. However, treating \Q@MyFoo@ as a closed rather than generic type is equivalent
to assuming that there is only one existential for all instances --- associated with the declaration site --- rather than potentially
different existentials for different instances. This combined with the fact that the JVM is unable to dynamically check for different instantiations
leads to the unsoundness.
Our dynamic solution resolves this with a specific implementation of casts on reified generics.
\mjnote{under modern Java there is no separate dynamic solution; see the note on the dynamic approach in the introduction.}
Our static solution effectively prohibits the de facto existential capture of the enclosing type parameter.


The decision to treat types like the \Q@MyFoo@ variants in our examples as closed existential types that depend on other choices of types is not
without parallels in other languages. Haskell and Rust, for example, include associated types~\cite{chakravarty2005associated}, where
a typeclass has a type as a member.
A yet closer analogue is Scala's path-dependent types~\tocite{dot}, where types may be \emph{members} of an individual object, rather than associated with an individual class. The code below adapts the troublesome example method using Scala's type members:
\begin{lstlisting}[language=Scala]
trait Foo { type X; def x() : X; }
def m[Y](other: Any, newx: Y) : Foo2 = 
    class MyFoo extends Foo { type X = Y; override def x() : X = newx; }
    newx match { 
        case f: MyFoo => f
        case _ => MyFoo()
    }
\end{lstlisting}
This code compiles without warnings --- and is not an issue. In this case the dependency of \Q@MyFoo@ instances
on the type parameter used to call \Q@m@ is made explicit with a type member. But Scala now treats thes

% Related work:
%  Big obvious task where you know much more then me, C#, Kotling etc
%
%  More related works:
%  Two other ways Java is unsound:
%  - The famous first paper
%
% -the well known snewaky throws
%
% public interface CheckedFunction<T, R> { R apply(T t) throws Exception; }
% @SuppressWarnings("unchecked")
% public static <T, E extends Exception> T sneakyThrow(Exception e) throws E {
%     throw (E) e;
% }
%
% public static <T, R> Function<T, R> unchecked(CheckedFunction<T, R> checkedFunction) {
%     return t -> {
%         try {
%             return checkedFunction.apply(t);
%         } catch (Exception e) {
%             return sneakyThrow(e);
%         }
%     };
% }
%
%
